% Part: normal-modal-logic
% Chapter: filtrations
% Section: euclidean-filtrations

\documentclass[../../../include/open-logic-section]{subfiles}

\begin{document}

\olfileid[tr]{nml}{fil}{euc}

\olsection{Öklidyen Modellerin Süzmeleri}

\olref[acc]{sec}'deki yaklaşım, Öklidyen ya da hem seri hem Öklidyen
modeller için işe yaramaz. \olref{fig:ser-eucl}'de gösterilen hem
Öklidyen hem seri modeli ele alalım. $\Gamma=\{p,\Box p\}$ olsun.
$\Gamma$ üzerinden süzme alındığında $w_1$ ile $w_3$, $\Gamma$ üzerinde
uyuşan tek farklı dünya çifti olduğundan $[w_1]=[w_3]$ olur. Her süzme,
$\mModel{M}$ modelinden devralınan ve \olref{fig:ser-eucl2}'de gösterilen
okları da içerir. Bu model Öklidyen değildir. Üstelik modeli Öklidyen yapmak
için ona oklar ekleyemeyiz. $[w_2]$ ile $[w_4]$ arasına iki yönlü oklar,
ardından $[w_2]$ ile~$[w_5]$ arasına da iki yönlü oklar eklememiz gerekirdi.
Fakat $\Box p$'nin $[w_2]$'de doğru, $p$'nin ise $[w_5]$'te yanlış olması
gerekir.

\begin{figure}[htpb]
  \centering
  \begin{tikzpicture}[modal]
    \node[world] (w1)
          [label={left: \mFalse{p}},
            label={below: $\mSat{{}}{\Box p}$}] {$w_1$};
    \node[world] (w2)
          [label={right: \mTrue{p}},
            label = {below: $\mSat{{}}{\Box p}$},
            right=of w1] {$w_2$};
    \draw[->] (w1) to (w2);
    \draw[reflexive above] (w2) to (w2);
    \node[world] (w3)
          [label={left: \mFalse{p}},
            label={below: $\mSat{{}}{\Box p}$},
            below=of w1] {$w_3$};
    \node[world] (w4)
          [label={right:\mTrue{p}},
            label={below: $\mSat/{{}}{\Box p}$},
            right=of w3] {$w_4$};
    \draw[->] (w3) to (w4);
    \draw[reflexive above] (w4) to (w4);
    \node[world] (w5)
          [label={right: \mFalse{p}},
            label=below:$\mSat/{{}}{\Box p}$,
            right=of w4] {$w_5$};
    \draw[->, bend left] (w4) to (w5);
    \draw[->, bend left] (w5) to (w4);
    \draw[reflexive above] (w5) to (w5);
  \end{tikzpicture}
  \caption{Seri ve Öklidyen bir model.}
  \ollabel{fig:ser-eucl}
\end{figure}

\begin{figure}[ht]
  \centering
  \begin{tikzpicture}[modal,every text node part/.style={align=center}]
    \node[world] (w1)
          [label={left: \mFalse{p}},
            label={right:$[w_1]=[w_3]$},
            label={below: $\mSat{{}}{\Box p}$}] {$[w_1]$};
    \node[world] (w2)
         [label={right: \mTrue{p}},
           label = {below: $\mSat{{}}{\Box p}$},
           right=of w1, yshift=1.5cm] {$[w_2]$};
    \draw[->] (w1) to (w2);
    \draw[reflexive above] (w2) to (w2);
    \node[world] (w4)
         [label={right:\mTrue{p}},
           label={below: $\mSat/{{}}{\Box p}$},
           right=of w1,yshift=-1.5cm] {$[w_4]$};
    \draw[->] (w1) to (w4);
    \draw[reflexive above] (w4) to (w4);
    \node[world] (w5)
         [label={right: \mFalse{p}},
             label=below:$\mSat/{{}}{\Box p}$,
             right=of w4] {$[w_5]$};
    \draw[->, bend left] (w4) to (w5);
    \draw[->, bend left] (w5) to (w4);
    \draw[reflexive above] (w5) to (w5);
  \end{tikzpicture}
  \caption{\olref{fig:ser-eucl}'deki modelin süzmesi.}
  \ollabel{fig:ser-eucl2}
\end{figure}

Özellikle Öklidyen bir süzme elde etmek için alt-!!{formula}s altında kapalı
gelişigüzel $\Gamma$ kümeleri üzerinden alınan süzmeleri incelemek yeterli
değildir. Bunun yerine \emph{modal olarak kapalı} $\Gamma$ kümelerini
incelememiz gerekir (bkz. \olref[pre]{defn:modallyclosed}). Bu tür tümce
kümeleri sonsuzdur; dolayısıyla karşılık gelen sistemin sonlu model
özelliğini ya da karar verilebilirliğini doğrudan vermezler.

\begin{thm}\ollabel{thm:modal-closed-filt}
  $\Gamma$ modal olarak kapalı,
  $\mModel{M}=\tuple{W,R,V}$ ve
  $\mModel{M^*}=\tuple{W^*,R^*,V^*}$, $\mModel{M}$'nin en kaba süzmesi
  olsun.
  \begin{enumerate}
  \item $\mModel{M}$ geçişliyse $\mModel{M^*}$ de geçişlidir.
  \item $\mModel{M}$ Öklidyense $\mModel{M^*}$ de Öklidyendir.
  \item $\mModel{M}$ simetrikse $\mModel{M^*}$ de simetriktir.
  \end{enumerate}
\end{thm}

\begin{proof}
\begin{enumerate}
  \item $\mModel{M^*}$ en kaba süzmeyse tanım gereğince $R^*[u][v]$
    ancak ve ancak $C_1(u,v)$ olduğunda geçerlidir. Geçişlilik için
    $C_1(u,v)$ ile $C_1(v,w)$'nin geçerli olduğunu varsayalım;
    $C_1(u,w)$'yi göstermeliyiz.
    \iftag{prvBox}{$\mSat{M}{\Box !A}[u]$ olsun. \Ax{4} bütün geçişli
      modellerde geçerli olduğundan $\mSat{M}{\Box\Box!A}[u]$; modal
      kapalılık nedeniyle $\Box\Box!A\in\Gamma$ olduğundan önce
      $C_1(u,v)$ ile $\mSat{M}{\Box!A}[v]$, sonra $C_1(v,w)$ ile
      $\mSat{M}{!A}[w]$ elde edilir. }{}
    \iftag{prvDiamond}{$\mSat{M}{!A}[w]$ olsun. Modal kapalılıktan
      $\Diamond!A\in\Gamma$ ve $C_1(v,w)$ ile
      $\mSat{M}{\Diamond!A}[v]$. Yine modal kapalılıktan
      $\Diamond\Diamond!A\in\Gamma$ ve $C_1(u,v)$ ile
      $\mSat{M}{\Diamond\Diamond!A}[u]$. $\Ax{4}_\Diamond$ bütün
      geçişli modellerde geçerli olduğundan
      $\mSat{M}{\Diamond!A}[u]$.}{}
  \item Alıştırma. Hem \Ax{5}'in hem de $\Ax{5_\Diamond}$'ın bütün
    Öklidyen modellerde geçerli olduğunu kullanın.
  \item Alıştırma. \Ax{B} ile $\Ax{B_\Diamond}$'ın bütün simetrik
    modellerde geçerli olduğunu kullanın.
\end{enumerate}
\end{proof}

\begin{prob}
  \olref[nml][fil][euc]{thm:modal-closed-filt} teoreminin ispatını
  tamamlayın.
\end{prob}

\end{document}
